\subsection{Capacidad de disparo sobre objetivo estacionario y en movimiento}

\subsubsection{Descripción de la prueba}

Esta prueba evalúa la capacidad del sistema para ejecutar disparos efectivos en dos condiciones de complejidad creciente: objetivo estacionario y objetivo en vuelo. Se busca verificar precisión, repetibilidad y tiempo de respuesta del conjunto visión--control--torreta--disparo.

\paragraph{Objetivos mínimos}
\begin{itemize}
    \item Impactar objetivo tipo dron estacionario en condiciones controladas.
    \item Impactar objetivo en movimiento con velocidad transversal representativa de operación.
    \item Verificar estabilidad del desempeño en secuencias múltiples de disparo.
\end{itemize}

\subsubsection{Criterios de aprobación}

\begin{table}[H]
\centering
\small
\begin{tabular}{p{8.8cm} p{4.1cm}}
\toprule
\textbf{Métrica} & \textbf{Criterio de aprobación} \\
\midrule
Tasa de impacto sobre objetivo estacionario (10 disparos) & $\geq 80\%$ \\
Tasa de impacto sobre objetivo en movimiento (10 pasadas) & $\geq 60\%$ \\
Tiempo de respuesta desde adquisición a primer disparo válido & $\leq 2{,}0\ \text{s}$ \\
Desviación angular media del punto de impacto respecto del centro del blanco & $\leq 4^\circ$ \\
Disparos fuera del sector de seguridad definido & 0 eventos \\
\bottomrule
\end{tabular}
\caption{Criterios de aprobación de la prueba de disparo.}
\label{tab:criterios_prueba_disparo}
\end{table}

\subsubsection{Procedimiento de ejecución}

\begin{enumerate}
    \item Delimitar campo de ensayo con sector de seguridad y zona de captura de impactos.
    \item Configurar el prototipo en modo de seguimiento y disparo asistido.
    \item Ejecutar una serie sobre objetivo estacionario y registrar impactos.
    \item Ejecutar pasadas de objetivo en movimiento, con trayectoria transversal controlada, y registrar impactos.
    \item Medir tiempos de adquisición y de disparo, junto con la desviación angular de impacto.
    \item Consolidar métricas para comparar contra los criterios de aprobación.
\end{enumerate}

\subsubsection{Resultados}

\begin{table}[H]
\centering
\small
\begin{tabular}{p{8.8cm} p{2.2cm} p{2.8cm}}
\toprule
\textbf{Métrica} & \textbf{Resultado} & \textbf{Cumple} \\
\midrule
Tasa de impacto objetivo estacionario [\%] & --- & --- \\
Tasa de impacto objetivo en movimiento [\%] & --- & --- \\
Tiempo de respuesta adquisición-disparo [s] & --- & --- \\
Desviación angular media de impacto [deg] & --- & --- \\
Eventos fuera de sector de seguridad [\#] & --- & --- \\
\bottomrule
\end{tabular}
\caption{Resultados medidos de la prueba de disparo.}
\label{tab:resultados_prueba_disparo}
\end{table}
